% =============================================================================
% Patch: §2 sec:local-kernel-and-jet-normalization
%
% Target file:   rh/rh_rebuild.tex
% Target lines:  307 through 471 (the entire current §2 block, including
%                §2.1 Local phase chart, §2.2 Phase kernel def, §2.3
%                Point-to-jet transform).
%
% Action: replace the existing §2 with the content below.
%
% Closes structurally:
%   - Identifies Phi = theta (Riemann-Siegel) on retained windows,
%     making the qualitative existence of the chart automatic.
%   - Adds lem:phase-kernel-properties (symmetry, real-valuedness,
%     removable-singularity).
%   - Adds lem:phase-kernel-derivatives (the kernel derivatives at
%     (T_1, T_2) and (T,T) used in §3).
%
% Does NOT close (residual under rem:wip-local-phase-chart):
%   - (P1) polynomial-in-Q lower bound on |zeta(1/2 + it)| on retained
%     packets — the chart-denominator condition.  Standard analytic
%     NT input; deferred.
%   - (P2) polynomial-in-Q lower bound on q = theta' on retained
%     packets — automatic at large T, but explicit constants deferred.
%
% Status changes:
%   \statusmig    not-started -> partial
%       (chart identification + kernel algebra closed; (P1)-(P2)
%        deferred to a later analytic-NT migration)
%   \statussympy  not-started -> verified
%       (verification: rh/sympy/sec-local-kernel-and-jet-normalization.py)
%   \statussim    not-started -> verified
%       (verification: rh/simulation/sec-local-kernel-and-jet-normalization.py)
%
% Verification scripts:
%   python rh/sympy/sec-local-kernel-and-jet-normalization.py
%   python rh/simulation/sec-local-kernel-and-jet-normalization.py
% =============================================================================

% =============================================================================
\section{Local kernel and jet normalization}
\label{sec:local-kernel-and-jet-normalization}
\statuspaper{LIVE}\quad
\statusmig{partial}\quad
\statusreferee{not-started}\quad
\statussympy{verified}\quad
\statussim{verified}\quad
\statuslean{not-started}
\par\medskip

A candidate off-line zero \(\rho\) at height \(\gamma\) is detected only
at scales comparable to \(1/\gamma\), and the analytic data of \(\zeta\)
beyond a window of width \(\asymp 1/Q(\gamma)\) around \(\gamma\) does not
enter the comparison object.  We therefore fix a real smooth phase
\(\Ph\) for \(\zeta\) on such a window once and for all, and never use
\(\zeta\) again --- only \(\Ph\), and only through the phase kernel
\(K_\Ph\) and its Taylor jets at one or two centers.

\subsection{Local phase chart}
\label{subsec:local-phase-chart}

\begin{definition}[Local phase chart]
\label{def:local-phase-chart}
For each height \(T\) we fix a real smooth function \(\Ph_T\) on an
interval \(I_T\) of length \(\asymp 1/Q(T)\) centered at \(T\) with the
following properties:
\begin{enumerate}[label=\textup{(P\arabic*)},leftmargin=*]
\item \(\Ph_T\) is the standard real critical-line phase of \(\zeta\) on
  \(I_T\), in a chart whose denominators admit a polynomial-in-\(Q\)
  lower bound;
\item the normalized derivative \(q:=\Ph_T'\) satisfies a
  polynomial-in-\(Q\) lower spectral bound on the retained packets at
  height \(T\).
\end{enumerate}
We write \(\Ph\) and \(q\) for \(\Ph_T,q\) when \(T\) is fixed by
context.
\end{definition}

\begin{remark}[Standard identification \(\Ph_T = \theta\)]
\label{rem:phase-chart-identification}
The local phase \(\Ph_T\) of Definition~\ref{def:local-phase-chart} is
the Riemann--Siegel theta function \(\theta\) restricted to \(I_T\).
The functional equation of \(\zeta\) gives the factorization
\[
\zeta(\tfrac12 + it) = e^{-i\theta(t)}\, Z(t),
\qquad
\theta(t) = \arg\Gamma\!\left(\tfrac14 + \tfrac{it}{2}\right) - \frac{t}{2}\log\pi,
\]
where \(Z\) is real-valued on \(\mathbb R\); we set \(\Ph_T(t) := \theta(t)\)
on \(I_T\).  The asymptotic expansion
\[
\theta(t) = \frac{t}{2}\log\frac{t}{2\pi} - \frac{t}{2} - \frac{\pi}{8}
          + \frac{1}{48 t} + O(t^{-3})
\]
gives in particular \(q(t) = \theta'(t) = \tfrac12 \log(t/(2\pi)) + O(t^{-2})\)
at large \(t\), so \(q\) is positive and grows logarithmically.
With \(\Ph_T=\theta\) the qualitative existence of the chart is
automatic; the residual content of (P1)--(P2) is the explicit
polynomial-in-\(Q\) constants.
\end{remark}

\begin{wip}\label{rem:wip-local-phase-chart}
With \(\Ph_T=\theta\) (Remark~\ref{rem:phase-chart-identification}),
the qualitative existence of the chart is automatic.  Two quantitative
inputs remain to be migrated:
\begin{enumerate}[label=\textup{(R\arabic*)},leftmargin=*]
\item (P1) requires a polynomial-in-\(Q\) lower bound on
  \(|\zeta(\tfrac12 + it)|\) on retained packets at height \(T\).
  This is the standard ``\(\zeta\) not too small'' condition; for a
  positive polynomial-weight subfamily of windows it follows from
  Selberg-type density estimates.  Quantitative migration deferred.
\item (P2) requires a polynomial-in-\(Q\) lower bound on
  \(q(t)=\theta'(t)\) on retained packets.  Since
  \(\theta'(t) = \tfrac12\log(t/(2\pi)) + O(t^{-2})\) at large \(t\),
  this is automatic for any \(Q\) polynomial in \(\log T\) and \(T\)
  sufficiently large.  Quantitative migration deferred.
\end{enumerate}
v1 prose to mine: the chart-construction parts of the local kernel
and jet normalization section, together with the curvature-estimate
material in the zeta-side decomposition section.
\end{wip}

\begin{definition}[Phase kernel]
\label{def:phase-kernel}
The phase kernel of \(\Ph\) on \(I_T\) is
\begin{equation}
\label{eq:phase-kernel}
K_\Ph(x,y)=
\begin{cases}
\dfrac{\sin\bigl(\Ph(x)-\Ph(y)\bigr)}{\pi(x-y)}, & x\ne y,\\[1.2ex]
\dfrac{q(x)}{\pi}, & x=y.
\end{cases}
\end{equation}
\end{definition}

\begin{lemma}[Phase-kernel basic properties]
\label{lem:phase-kernel-properties}
Let \(\Ph\in C^3(I_T;\mathbb R)\).  The phase kernel \(K_\Ph\) of
Definition~\ref{def:phase-kernel} satisfies:
\begin{enumerate}[label=\textup{(\roman*)},leftmargin=*]
\item \emph{symmetry}: \(K_\Ph(x,y) = K_\Ph(y,x)\);
\item \emph{real-valuedness}: \(K_\Ph(x,y)\in\mathbb R\) for all
  \((x,y)\in I_T\times I_T\);
\item \emph{removable singularity}:
  \(\lim_{x\to y} K_\Ph(x,y) = q(y)/\pi\), and the displayed
  diagonal value \eqref{eq:phase-kernel} is the unique continuous
  extension across the diagonal.
\end{enumerate}
\end{lemma}

\begin{proof}
(i) For \(x\ne y\),
\[
K_\Ph(y,x) = \frac{\sin(\Ph(y)-\Ph(x))}{\pi(y-x)}
           = \frac{-\sin(\Ph(x)-\Ph(y))}{-\pi(x-y)}
           = K_\Ph(x,y).
\]
(ii) The numerator \(\sin(\Ph(x)-\Ph(y))\) is real because \(\Ph\) is
real-valued; the denominator \(\pi(x-y)\) is real; the diagonal value
\(q(y)/\pi\) is real.  Hence \(K_\Ph\) is real-valued throughout.
(iii) Set \(x = y+h\) with \(h\to 0\).  Taylor-expanding \(\Ph\)
around \(y\),
\[
\Ph(y+h) - \Ph(y) = h\,q(y) + \tfrac{h^2}{2}q'(y) + O(h^3),
\]
so
\[
\frac{\sin(\Ph(y+h)-\Ph(y))}{\pi h}
= \frac{h\,q(y) + O(h^2)}{\pi h}
= \frac{q(y)}{\pi} + O(h),
\]
which converges to \(q(y)/\pi\) as \(h\to 0\).  Continuity of \(\Ph\)
and \(q\) makes the displayed extension the unique continuous one.
\end{proof}

\begin{lemma}[Phase-kernel derivatives at distinct points]
\label{lem:phase-kernel-derivatives}
At distinct points \(T_1 \ne T_2\), with \(s = T_1 - T_2\),
\(\Del = \Ph(T_1) - \Ph(T_2)\), and \(q_j = q(T_j)\), the partial
derivatives of \(K_\Ph\) at \((T_1, T_2)\) are
\begin{align*}
K_\Ph(T_1, T_2)
  &= \frac{\sin\Del}{\pi s},\\
\partial_x K_\Ph(T_1, T_2)
  &= \frac{q_1\, s\cos\Del - \sin\Del}{\pi s^2},\\
\partial_y K_\Ph(T_1, T_2)
  &= \frac{\sin\Del - q_2\, s\cos\Del}{\pi s^2},\\
\partial_{xy} K_\Ph(T_1, T_2)
  &= \frac{(q_1+q_2)s\cos\Del + (q_1 q_2 s^2 - 2)\sin\Del}{\pi s^3}.
\end{align*}
\end{lemma}

\begin{proof}
Write \(K_\Ph(x,y) = N(x,y)/D(x,y)\) with \(N(x,y) = \sin(\Ph(x)-\Ph(y))\)
and \(D(x,y) = \pi(x-y)\).  Apply the quotient rule and the chain rule
on \(\sin\), using \(\partial_x \Ph(x) = q(x)\) and
\(\partial_y \Ph(y) = q(y)\).  The first three identities are direct;
the mixed second derivative \(\partial_{xy}\) follows by differentiating
\(\partial_x K_\Ph\) in \(y\), keeping the four terms (carrier \(K\),
factor \(\sin\Del\), factor \(\cos\Del\), denominator power) and
collecting.  Symbolic verification entry by entry is in
\texttt{rh/sympy/sec-local-kernel-and-jet-normalization.py}.
\end{proof}

\begin{remark}[Diagonal limit of the kernel derivatives]
\label{rem:kernel-derivatives-diagonal}
As \(T_2 \to T_1 = T\), the formulas of
Lemma~\ref{lem:phase-kernel-derivatives} take the diagonal limits
\[
K_\Ph(T,T) = \frac{q(T)}{\pi},\qquad
\partial_x K_\Ph(T,T) = \partial_y K_\Ph(T,T) = \frac{q'(T)}{2\pi},
\qquad
\partial_{xy} K_\Ph(T,T) = \frac{q''(T) + 2q(T)^3}{6\pi},
\]
which is the kernel-derivative form of \(J(T)\) used in
\eqref{eq:same-point-as-kernel-derivatives}.  These limits are
uniform on retained windows by (P1)--(P2) of
Definition~\ref{def:local-phase-chart}.
\end{remark}

\subsection{Point-to-jet transform}
\label{subsec:point-to-jet-transform}

The kernel \(K_\Ph\) is used only at small symmetric pairs \(T\pm h\).
We rotate the \(2\times 2\) block at such a pair out of the point basis
by the \emph{Gram-normalized coalescing even/odd transform}
\begin{equation}
\label{eq:point-to-jet-transform}
P_h
\;=\;
\frac1{\sqrt2}
\begin{pmatrix}
1 & 1\\[0.6ex]
-1/(2h) & 1/(2h)
\end{pmatrix}.
\end{equation}
Equivalently, \eqref{eq:point-to-jet-transform} is obtained from the
literal central value/derivative jet transform by a fixed diagonal
congruence:
\begin{equation}
\label{eq:gram-vs-jet-congruence}
P_h
\;=\;
\diag(\sqrt 2,\, 1/\sqrt 2)\,
\begin{pmatrix}
1/2 & 1/2\\[0.6ex]
-1/(2h) & 1/(2h)
\end{pmatrix}.
\end{equation}
With this convention the same-point limit is
\begin{equation}
\label{eq:same-point-as-kernel-derivatives}
J(T)
\;=\;
\begin{pmatrix}
2K_\Ph(T,T) & \partial_y K_\Ph(T,T)\\[0.6ex]
\partial_x K_\Ph(T,T) & \tfrac12 \partial_{xy} K_\Ph(T,T)
\end{pmatrix},
\end{equation}
which evaluates entry by entry to the explicit phase-kernel formula in
Lemma~\ref{lem:same-point-jet-limit}.  We refer to the basis produced
by \(P_h\) as the \emph{Gram-normalized basis}, not the literal
value/derivative basis: it differs from the latter by the diagonal
congruence \eqref{eq:gram-vs-jet-congruence}, which carries the
factor-of-\(2\) rescaling visible in the leading entry of \(J(T)\) and
the corresponding rescalings of every cross-block and finite-scale
block downstream.

\begin{remark}[Pair separation convention]
\label{rem:pair-separation-convention}
We use the symmetric pair \(T\pm h\), not \(T\pm h/2\); the pair
separation is therefore \(2h\), and the \(1/(2h)\) factor in the lower
row of \(P_h\) is the inverse separation.  All formulas in
Sections~\ref{sec:jet-limit-local-blocks}
and~\ref{sec:finite-scale-local-blocks-and-whitening} are stated under
that convention.  Any imported lemma using the alternative convention
\(T\pm h/2\) (or using a \(P_h\) with \(1/h\) in the lower row) must be
retransformed before insertion.
\end{remark}

\begin{remark}[Convention audit: Gram-normalized vs.\ literal jet]
\label{rem:normalization-audit}
The downstream object of this paper is a whitened Gram block
\(\widehat\Om_\z(s;m)\)
(Definition~\ref{def:whitened-finite-block}); Gram-normalizing the
even and odd point modes before block extraction is the cleaner
Hilbert-space convention.  The displayed jet-limit blocks --- the
\(2 q(T)/\pi\) leading entry and the \(\tfrac1{12}(q''+2q^3)/\pi\)
corner of \(J(T)\) (Lemma~\ref{lem:same-point-jet-limit}), and the
matching cross-block entries \(2\sin\Del/(\pi s)\) etc.\ of
\(N_{12}\) (Lemma~\ref{lem:cross-block-jet-limit}) --- are the
Gram-normalized blocks; they are the limits of \(P_h A_h P_h^\top\)
for \eqref{eq:point-to-jet-transform}, not for the literal-jet
alternative.

The earlier draft \texttt{rh/proof\_of\_rh.tex} contains a convention
collision at exactly this point: it prints the literal-jet
transform
\(P_h^{\mathrm{old}}=\tfrac12\bigl(\begin{smallmatrix}1&1\\-1/h&1/h\end{smallmatrix}\bigr)\)
and immediately afterward displays the Gram-normalized \(J(T)\) used
above.  Those two choices do not match.  With \(P_h^{\mathrm{old}}\)
and the pair \(T\pm h\), conjugation of the raw kernel block produces
\[
\frac1\pi
\begin{pmatrix}
q(T) & \tfrac12 q'(T)\\[0.6ex]
\tfrac12 q'(T) & \tfrac16\bigl(q''(T)+2q(T)^3\bigr)
\end{pmatrix},
\]
which is the literal central value/derivative jet block, \emph{not}
the displayed \(J(T)\).  The two differ by the diagonal congruence
\eqref{eq:gram-vs-jet-congruence}: rescale the \((1,1)\) entry by
\(2\) and the \((2,2)\) entry by \(1/2\).  Not a typo in either
matrix --- a convention collision.

The present rebuild fixes the collision in favor of the Gram-normalized
transform.  All same-point and cross-block jet limits, all finite-scale
blocks \(G_{m,\pm}(s)\) and \(N_m(s)\), the whitened block
\(\widehat\Om_\z(s;m)\), and every quotient covector are stated under
\eqref{eq:point-to-jet-transform}.  The literal-jet alternative would
also be valid, but only after rescaling each of these objects
consistently; that rescaling is not done here.  Any imported v1 lemma
or formula that silently assumes \(P_h^{\mathrm{old}}\) must be
re-conjugated by \(\diag(\sqrt 2,\,1/\sqrt 2)\) before insertion.
\end{remark}
